\section{Purpose of the work}

The aim of this laboratory work is to measure how the number of threads and the way the work is
divided among them change the running time and the memory of a real workload: counting the
negative, positive and zero values among 50\,000\,000 integers stored in a text file. From my own
measurements I decide when this work is CPU-bound and when it is I/O-bound, I observe a race
condition and the cost of fixing it, and I find where the parallel speed-up stops.

\section{Source code}

The program is written in Java and is kept in the Git repository
\begin{center}
\url{https://github.com/bogdan-arama/pr-labs} \quad (folder \texttt{lab1})
\end{center}
It is one program, \texttt{lab1.Main}; every task and configuration is chosen by command-line
parameters, for example \texttt{java -cp out lab1.Main 1a --threads 8}. The
\texttt{README.md} in the folder explains how to build it (\texttt{javac -d out src/lab1/*.java}),
how to generate the data and how to run every task. Two small Python scripts run all
configurations (\texttt{scripts/run\_benchmarks.py}) and draw the charts of this report
(\texttt{scripts/plot.py}); all raw measurements are in \texttt{results/*.csv}.

\section{Work environment}

\begin{table}[H]
\centering\small
\begin{tabular}{@{}ll@{}}
\toprule
Item & Value \\
\midrule
Computer & laptop, connected to mains power \\
CPU & AMD Ryzen 5 5600H (Zen 3), 3.3\,GHz base, up to 4.2\,GHz boost \\
Cores & 6 physical cores, 12 logical processors (SMT, 2 threads per core) \\
RAM & 16\,GB DDR4-3200 (2 modules, 15.4\,GB usable) \\
Disk & Intel SSD 670p 512\,GB (\texttt{SSDPEKNU512GZ}), NVMe, PCIe 3.0 x4 \\
Operating system & Windows 11 Pro, build 10.0.22624, 64-bit \\
Language and runtime & Java 17, Eclipse Temurin OpenJDK 17.0.18+8, HotSpot 64-Bit Server VM \\
Helper scripts & Python 3.14.2, matplotlib 3.10.8 \\
\bottomrule
\end{tabular}
\caption{Hardware and software used for all measurements}
\end{table}

\section{Data and method of measurement}

\textbf{Data.} \texttt{java -cp out lab1.Main gen} writes \texttt{data/numbers.txt}: 50\,000\,000
integers drawn uniformly from $[-1\,000\,000,\,1\,000\,000]$ with \texttt{SplittableRandom}
(seed 42), one per line, 369\,445\,634 bytes. While writing, the generator also counts the signs
and saves them in \texttt{numbers.txt.expected}.

\textbf{Correctness check.} The expected result is
\textbf{25\,001\,107 negative, 24\,998\,868 positive and 25 zero} values, which add up to
50\,000\,000. After every run the program compares its totals with this file and prints
\texttt{correct=true} or \texttt{false}; the benchmark script stops if any correct version gives
wrong totals. Every run of Tasks 1a, 1b (versions 2 and 3) and 1c produced exactly these
numbers.

\textbf{Time.} \texttt{System.nanoTime()} is read just before the threads are created and just
after the last one has been joined. JVM start-up is not included, and in Task 1b loading the array
is not included either, as the task requires.

\textbf{Peak memory.} At the end of every run the program asks the operating system for the peak
working set of its own process (\texttt{PeakWorkingSet64} on Windows, \texttt{VmHWM} on Linux).
This is the largest amount of physical memory the whole process (JVM, heap, thread stacks, native
buffers) used during the run. To make it belong to one configuration only, every configuration is
run in a fresh JVM.

\textbf{Repetitions.} Tasks 1a and 1c were run 5 times per thread count, Task 1b 3 times; the charts show the median and a shaded band from the fastest to the
slowest run. Repetitions were interleaved (all thread counts once, then all again), so slow
changes of the machine (temperature, background programs) spread over all configurations. Besides
the thread counts required by the task (1, 2, 4, 8, 16, 32, 64) I also measured 6 and 12, the
numbers of physical cores and logical processors. In the charts these two are marked with dotted
vertical lines. Speed-up is the median time on 1 thread divided by the median time on $N$ threads.

\textbf{File cache.} The 369\,MB file easily fits in the 16\,GB of RAM, so after the first read
Windows keeps it in its file cache. Before Tasks 1a and 1c the benchmark makes one unmeasured run
to bring the file into the cache, so all measured runs read it from memory, not from the SSD.
